% ============================================================
% sec/01contexto.tex — Contexto (aprox. 2 min)
% ============================================================
\section{Contexto}

\begin{frame}{?`Qu\'e es el reconocimiento facial?}
  \framesubtitle{Un rostro se convierte en un vector y la decisi\'on es un umbral}
  \decidir{Expone: nombre \textbullet\ aprox. 2 min toda la secci\'on}

  \vspace{0.2cm}
  \begin{center}
  \begin{tikzpicture}[
      paso/.style={draw=tecsteel, fill=boxbg, rounded corners=2pt,
                   text width=1.75cm, minimum height=1cm, align=center,
                   font=\scriptsize},
      clave/.style={paso, fill=tecnavy, text=white, draw=tecnavy},
      flecha/.style={-{Stealth[length=2mm]}, tecsteel, thick},
      node distance=0.35cm]
    \node[paso] (a) {Captura};
    \node[paso, right=of a] (b) {Detecci\'on del rostro};
    \node[paso, right=of b] (c) {Alineaci\'on};
    \node[paso, right=of c] (d) {Representaci\'on (vector)};
    \node[clave, right=of d] (e) {Distancia vs.\ umbral};
    \node[paso, right=of e] (f) {Decisi\'on};
    \foreach \x/\y in {a/b,b/c,c/d,d/e,e/f} \draw[flecha] (\x) -- (\y);
  \end{tikzpicture}
  \end{center}

  \vspace{0.3cm}
  \begin{columns}[T]
    \begin{column}{0.46\textwidth}
      \begin{block}{Verificaci\'on (1:1)}
        ?`Esta persona es quien dice ser?\\
        {\scriptsize\color{labelgray}Desbloqueo del tel\'efono, control
        fronterizo.}
      \end{block}
    \end{column}
    \begin{column}{0.46\textwidth}
      \begin{block}{Identificaci\'on (1:N)}
        ?`Qui\'en es esta persona?\\
        {\scriptsize\color{labelgray}B\'usqueda contra listas de vigilancia
        en espacio p\'ublico.}
      \end{block}
    \end{column}
  \end{columns}
\end{frame}

\begin{frame}{?`Por qu\'e importa hoy?}
  \begin{columns}[T]
    \begin{column}{0.52\textwidth}
      \begin{itemize}\setlength{\itemsep}{6pt}
        \item Del tel\'efono al espacio p\'ublico: c\'amaras que identifican
              sin interacci\'on ni consentimiento.
        \item Un rasgo biom\'etrico no se revoca como una contrase\~na.
        \item El error no es abstracto: un falso positivo recae sobre una
              persona concreta.
        \item Regulaci\'on dispar: la UE restringe la identificaci\'on
              biom\'etrica remota; Costa Rica se rige por la Ley 8968
              (2011).
      \end{itemize}
    \end{column}
    \begin{column}{0.42\textwidth}
      \begin{block}{Tr\'iada CID}
        \scriptsize
        \textbf{C} --- Una base de rostros expuesta filtra un dato
        irrevocable.\\[3pt]
        \textbf{I} --- Una falsa coincidencia o una suplantaci\'on alteran
        la identidad atribuida.\\[3pt]
        \textbf{D} --- Un falso rechazo niega el acceso a un servicio.
      \end{block}
    \end{column}
  \end{columns}
\end{frame}